\chapter{Marco teórico}\label{cap3}
Las personas suelen pensar en la electricidad como algo que fluye. De hecho, eso es realmente cierto ,Relacionando la energía del sol como una fuente de energía, se obtiene que el sol emite radiación electromagnética que viaja como lo hacen las olas en el agua [1] ,  también hay que aclarar que solamente una cantidad de la energía solar que se emite incide parcialmente sobre la superficie de la  tierra y aun cuando llega a la incidir sobre la superficie terrestre, es rebotada de vuelta hacia el espacio, otra es absorbida por el aire, pero la mayoría del tiempo la energía solar que incide sobre la tierra es luz visible, de la cual se puede sacar provecho transformándola  en energía eléctrica utilizable.\\

La luz visible es transformada en energía eléctrica mediante paneles solares, pero no siempre esta transformación es efectuada de la mejor manera, dado a que hay factores que hacen que la captación solar no genere toda la energía que es capaz un panel solar de generar. El polvo es uno de los principales agentes que afectan la captación de luz, Experimentos realizados en parques solares, por ejemplo, muestran que la eficiencia de un panel fotovoltaico puede decrecer un 40\% por la acumulación de residuos en la superficie [5].\\

Otro estudio fue conducido para obtener datos acerca de como el polvo depositado en las superficies de paneles solares ha degradado el desempeño en sistemas solares, el objetivo era obtener datos sobre cuanto polvo podía incidir sobre la superficie del panel y la condiciones del ambiente dando como resultado que a diferencia de los paneles que contienen un nivel alto de obstrucción por el polvo, los paneles solares que obtienen una incidencia de luz solar propicia tienden a generar más energía y a tener una temperatura más alta [11].\\

Debido al incremento paulatino que se ha dado en el uso de la energía eléctrica convencional, se ha visto necesario crear fuentes alternativas que tengan ventajas, principalmente ventajas de las cuales se tenga en cuenta los beneficios para el medio ambiente, el uso de energías renovables en sistemas ha atraído bastante atención en los últimos años. En especial uno de los artefactos que más han atraído la atención debido a sus propiedades en pro del ambiente y de la generación limpia de energía eléctrica han sido las celdas fotovoltaicas, remarcando su capacidad de aprovechar de una manera especial la luz solar para convertirla en una utilidad para todos en forma de electricidad.\\

En los sistemas fotovoltaicos, los elementos conocidos como celdas fotovoltaicas son posicionadas en paneles expuestos a la luz solar y transforman la energía a través del flujo de electrones entre dos capas de semiconductores [2].\\
El proceso de obtener electricidad desde las celdas solares empieza cuando los átomos del silicón absorben la luz, la energía de la luz mueve algunos electrones fuera de los átomos, lo que hace que haya una fluidez entre los materiales semiconductores del panel solar, esta fluidez crea una corriente eléctrica que puede abandonar la celda mediante un material metálico pero la fluidez se ve interrumpida o cortada por la suciedad acumulada sobre los paneles haciendo que en veces el panel no trabaje al máximo de su capacidad.\\  
%  \ref{fig1}.
\begin{savenotes}
\begin{figure}[H]
	\centering
	\includegraphics[width=0.7\textwidth,keepaspectratio]{celda_solar}
	\caption{Partes de una celda solar.}
	\label{fig1}
\end{figure}
\end{savenotes}
\section{Paneles solares}
Una simple celda solar crea solo un poco de energía aprovechable, para la mayoría de los casos es necesaria más energía de la que ésta puede proveer, por lo cual las celdas están conectadas en grupos, estos grupos son llamados módulos solares.\\

Estos módulos solares son de pocos metros de largo y ancho y pueden producir un tanto de miles de watts de electricidad, pero una ventaja que explotan estos módulos es que pueden ser juntados para crear un arreglo solar que pueda producir más energía.
\begin{savenotes}
\begin{figure}[H]
	\centering
	\includegraphics[width=0.7\textwidth,keepaspectratio]{celda_solar_2}
	\caption{Arreglos de celdas solares.}
	\label{fig2}
\end{figure}
\end{savenotes}
Ahora bien, los paneles solares se instalan de manera en que la incidencia de la luz solar sea aprovechada de la manera más eficiente para generar la mayor cantidad de energía posible, pero no solo la posición del panel es importante, también hay agentes externos que entorpecen la producción de electricidad.\\

La energía o eficiencia producida por módulos de celdas fotovoltaicas es relacionado con la disponibilidad de la irradiación solar, así como otros factores ambientales, clima, desempeño de componentes y sistemas inherentes. Este polvo, tierra y heces de aves son las razones con más mérito por la obstrucción del desempeño de sistemas fotovoltaicos [3].\\

El polvo cubre paneles solares y elementos ópticos, tales como lentes y espejos, causando la degradación de su desempeño óptico [4].
\subsection*{Proceso de Limpieza}
EL proceso de limpieza de un panel solar se centra en tener en cuenta que un panel solar puede llegar ser un elemento frágil propenso a que si usamos ciertos tipos de agentes limpiadores el cristal de la parte superior del panel se vea afectada y desgastada, además de que debe ser lo suficientemente eficaz para limpiar fluidos que al secarse son más difíciles de limpiar.\\

Experimentos realizados en parques solares, por ejemplo, muestran que la eficiencia de un panel fotovoltaico puede decrecer un 40\% por la acumulación de residuos en la superficie [5].\\

Por lo tanto, para que las celdas fotovoltaicas operen a su máxima eficiencia sin pérdidas de energía, es necesario que los fotones tengan un libre acceso al panel de celdas fotovoltaicas [5].

\section{Sistemas}
La Mecatrónica trabaja con lo que se conoce como sistema. Un sistema puede considerarse una caja con una entrada y una salida y de la cual no nos enteramos de su contenido, sino la relación que existe entre la salida y la entrada.\\

Un \textbf{sistema de medición} se podría considerar como una caja negra que se utiliza para medir. Su entrada es la magnitud que se desea medir y su salida es el valor correspondiente a dicha magnitud.  En general puede decirse que los sistemas de medición están formados por tres elementos:
\begin{enumerate}
    \item Un sensor.
    \item Un acondicionador de señal.
    \item Un sistema de presentación visual.
\end{enumerate}
Un \textbf{sistema de control} puede considerarse como una caja negra que sirve para controlar la salida de un valor o secuencia de valores determinados. Existen dos tipos de sistemas de control: el de lazo abierto y el de lazo cerrado. Los sistemas de lazo abierto tienen la ventaja de ser relativamente sencillos, por lo que su costo es bajo y en general su confiabilidad es buena. Sin embargo, con frecuencia son imprecisos ya que no hay corrección de errores. Los sistemas en lazo cerrado tienen la ventaja de ser bastante precisos para igualar el valor real y el deseado. Pero son más complejos y, por lo tanto, más costosos y con mayor probabilidad de descompostura debido a la mayor cantidad de componentes.
\subsection{Sistemas mecánicos}
Los mecanismos, es decir, dispositivos que se pueden considerar de movimiento, en tanto transforman el movimiento de una forma a otra, con un mecanismo un movimiento lineal se puede convertir en un movimiento rotacional; un movimiento que se produce en una dirección en uno con otra dirección en ángulo recto respecto de la primera; un movimiento lineal alterno en uno rotacional. Entre los elementos mecánicos están los de barras articuladas, levas, engranes, cremalleras, cadenas, correas de transmisión, etc.
\subsection{Sistemas eléctricos}
Al estudiar los sistemas eléctricos que se emplean como actuadores de control deberán tenerse en cuanta los siguientes dispositivos y sistemas:
\begin{enumerate}
    \item Dispositivos de conmutación, como son los interruptores mecánicos (relevadores) y los interruptores de estado sólido (diodos, tiristores y transistores), en los que la señal de control enciende o apaga un dispositivo eléctrico.
    \item Dispositivos tipo solenoide, en los cuales una corriente que pasa por un solenoide acciona un núcleo de hierro dulce.
    \item Sistemas motrices, como motores de CD y de CA, en los cuales la corriente que pasa por el motor produce una rotación.
\end{enumerate}
\section{Automatización}
Para que los componentes necesarios para el funcionamiento de una máquina se integren en una perfecta armonía, hace falta más que integrarlos físicamente al mecanismo es necesario lograr el control de estos para que la máquina se comporte de la manera deseada.\\

El control automático ha desempeñado una función vital en el avance de la ingeniería y la ciencia. El control automático se ha vuelto una parte importante e integral de los procesos modernos industriales y de manufactura.\\

Debido a que los avances en la teoría y la práctica del control automático aportan los medios para obtener un desempeño óptimo de los sistemas dinámicos, mejorar la productividad, aligerar la carga de muchas operaciones manuales repetitivas y rutinarias, así como de otras actividades, casi todos los ingenieros y científicos deben tener un buen conocimiento de este campo.\\

El primer trabajo significativo en control automático fue el regulador de velocidad centrífugo de James Watt para el control de la velocidad de una máquina de vapor. En 1922, Minorsky trabajó en los controladores automáticos para dirigir embarcaciones, y mostró que la estabilidad puede determinarse a partir de las ecuaciones diferenciales que describen el sistema.\\

En 1932, Nyquist diseñó un procedimiento relativamente simple para determinar la estabilidad de sistemas en lazo cerrado, con base en la respuesta en lazo abierto en estado estable cuando la entrada aplicada es una senoidal. En 1934, Hazen, quien introdujo el término servomecanismos para los sistemas de control de posición, analizó el diseño de los servomecanismos con relevadores, capaces de seguir con precisión una entrada cambiante.\\

Durante la década de los cuarenta, los métodos de la respuesta en frecuencia hicieron posible que los ‘ingenieros diseñan sistemas de control lineales en lazo cerrado que cumplieran con los requerimientos de desempeño. A finales de los años cuarenta y principios de los cincuenta, ‘se desarrolló por completo el método del lugar geométrico de las raíces propuesto por Evans.\\

Los métodos de respuesta en frecuencia y del lugar geométrico de las raíces, que forman el núcleo de la teoría de control clásica, conducen a sistemas estables que satisfacen un conjunto más o menos arbitrario de requerimientos de desempeño. En general, estos sistemas son aceptables, pero no óptimos en forma significativa.
Conforme las plantas modernas con muchas entradas y salidas se vuelven más y más complejas, la descripción de un sistema de control moderno requiere de una gran cantidad de ecuaciones. La teoría del control clásica, que trata de los sistemas con una entrada y una salida, pierde su solidez ante sistemas con entradas y salidas múltiples Desde alrededor de 1960, debido a que la disponibilidad de las computadoras digitales hizo posible el análisis en el dominio del tiempo de sistemas complejos, la teoría de control moderna, basada en el análisis en el dominio del tiempo y la síntesis a partir de variables de estados, se ha desarrollado para enfrentar la creciente complejidad de las plantas modernas y los requerimientos limitativos respecto de la precisión, el peso y el costo en aplicaciones militares, espaciales e industriales.\\

Durante los años comprendidos entre 1960 y 1980, se investigaron a fondo el control óptimo tanto de sistemas determinísticos como estocásticos, y el control adaptable, mediante el aprendizaje de sistemas complejos. De 1980 a la fecha, los descubrimientos en la teoría de control moderna se centraron en el control robusto, el control de H, y temas asociados.\\

Ahora que las computadoras digitales se han vuelto más baratas y más compactas, se usan como parte integral de los sistemas de control. Las aplicaciones recientes de la teoría de control moderna incluyen sistemas ajenos a la ingeniería, como los biológicos, biomédicos, económicos y socioeconómicos.

\section{Métodos de limpieza de paneles solares}
Numerosos tipos de métodos de limpieza están disponibles en el mercado. Cada uno de estos tiene ventajas y desventajas. Esta sección proporciona una breve explicación de las posibles alternativas de limpieza.
\subsection{Mecánico}
La limpieza mecánica es el tipo más común, que incluye robots, mano de obra, cepillos, etc. normalmente necesita un operador. Estos robots tienen cepillo integrado, que no necesita mucha agua y son pequeños. Las máquinas autónomas son robustas; sin embargo, necesitan personalización y mucha mano de obra. Otro tipo de robots son los que están basados en rociadores, que necesitan más agua para limpiar los paneles. Mientras que la limpieza manual, es generalmente utilizada para limpiar paneles solares domésticos o pequeñas instalaciones. Inaccesibilidad al agua y los cambios en los costos de mano de obra, así como las condiciones climáticas extremas, pueden considerarse como las principales desventajas de este método. Sin embargo, las versiones automatizadas de los métodos de limpieza mecánica pueden ser más flexibles con estos.
\subsection{Eléctrico}
El método de eliminación de polvo por vía eléctrica más popular se basa en el "concepto de cortina eléctrica" desarrollado por F.B. Tatom y sus colaboradores en la NASA en 1967.
El concepto básico de esta solución es aplicar una onda eléctrica que forma una onda en movimiento que bloquearía las partículas que se asientan en el módulo fotovoltaico [39]. Una desventaja es que la superficie del panel se cargará en todo momento, lo que no sería seguro durante los días lluviosos.

\subsection{Autolimpieza química}
Las nano películas autolimpiables pueden estar en forma de superhidrofobicidad y películas superhidrofóbicas. La durabilidad y la estabilidad de este tipo de materiales de recubrimiento son parámetros importantes para operaciones a largo plazo. Primero, el proceso fotocatalítico tendría lugar, donde la película reacciona bajo luz ultravioleta y divide la suciedad orgánica. Después de eso, el agua de lluvia elimina el polvo. Sin embargo, esta solución requiere lluvia para enjuagar el polvo de los paneles solares, lo que significa que necesita agua en regiones secas. Otra desventaja es una pérdida de eficiencia de aproximadamente 0.2-0.3% durante el período de dos meses.
\subsection{Limpieza por succión al vacio}
Una aspiradora de vacío es un dispositivo que utiliza una bomba de aire para crear un vacío parcial para aspirar el polvo y la suciedad, generalmente de pisos, ventanas, etc. En general, el suministro eléctrico se proporciona al motor de la aspiradora que crea la presión de succión. La desventaja es el consumo de energía de la aspiradora es en watts y no justifica la efectividad de la aspiradora, aunque cumple con el objetivo el gasto de energía es muy alto.
La potencia de entrada se convierte en flujo de aire al final y se mide en watts de aire. Además, la aspiradora puede limpiar el panel correctamente solo en las superficies que no sean las esquinas y esto debe manejarse manualmente. Durante un período de tiempo, los rasguños y el polvo acumulado causan una absorción deficiente de la incidencia de radiación solar.
\subsection{Limpieza automática basada en limpiaparabrisas}
La limpieza automática a base de limpiaparabrisas incorpora un limpiador de goma y un recipiente para rociar agua con aditivos. El proceso es exactamente como la limpieza del vidrio de un carro y requiere un mecanismo automático para operar y completar la tarea. El mecanismo funciona con pilas. Este método es similar al anterior y funciona automáticamente mediante el mecanismo de control adecuado, pero los impactos son similares a los anteriores.
\subsection{Precipitador electrostático}
Teniendo en cuenta los impactos de diferentes métodos de limpieza para paneles solares fotovoltaicos, los daños mecánicos, arañazos y cicatrices como las marcas del precipitador electrostático funcionan con un mecanismo sin contacto. Limpia eficientemente y protege la superficie superior del panel sin ningún contacto físico. Un precipitador electrostático (ESP) es un dispositivo de filtración que elimina partículas finas de polvo de la superficie del panel solar fotovoltaico utilizando la fuerza de una carga electrostática inducida. El mecanismo consiste en los electrodos.
\section{Sensores}
Un sensor es un dispositivo que, a partir de la energía del medio donde se mide, entrega una señal de salida transducible que es función de la variable medida. Sensor y transductor se emplean a veces como sinónimos, pero “sensor” sugiere un significado más extenso: la ampliación de los sentidos para adquirir un conocimiento de cantidades físicas que, por su naturaleza o tamaño, no pueden ser percibidas directamente por los sentidos.
\subsection{RTD}
Los detectores de temperatura basados en la variación de una resistencia eléctrica se suelen designar con sus siglas en ingles RTD (Resistance Temperature Detector).
El empleo de un conductor para la medida de temperaturas, está sometido a varias limitaciones. En primer lugar, es obvio que no se podrán medir temperaturas próximas ni superiores a la de fusión del conductor. En segundo lugar, para poder medir una temperatura determinada con este método, es necesario que el sensor este a dicha temperatura precisamente. Otra limitación a considerar es la posible presencia de deformaciones mecánicas, que, como se ha visto al exponer las galgas extensiométricas, provocan también un cambio en el valor de la resistencia eléctrica de un conductor.
\subsection{LDR}
Las fotorresistencias o fotoconductores (en inglés, Light Dependent Resistors, LDR) se basan en la variación de la resistencia eléctrica de un semiconductor al incidir en la radiación óptica.
La energía necesaria para producir el salto puede venir de otras fuentes externas además del calor, como pueden ser una radiación óptica o una tensión eléctrica. Los fotoconductores más comunes, utilizables a temperatura ambiente, son el SCd, SPb y SePb (en particular el primero). Se fabrican en formas muy variadas, con dispositivos simétricos, diferenciales, entre otros, para facilitar su aplicación. Las aplicaciones de las LDR ordinarias se pueden dividir en las de medición de luz (con poca precisión y bajo coste) y las que emplean la radiación para modificar señales. En el primer grupo están el control automático de brillo y contraste en receptores de televisión, control de diafragmas de cámaras fotográficas, detección de fuego, control de iluminación de vías públicas, etc. En el segundo grupo están los detectores de presencia y posición, y algunas mediciones de nivel de depósitos.
\subsection{Sensor de Polvo}
Este sensor de Sharp (GP2Y1010AU0F) es especialmente efectivo detectando partículas muy finas de polvo en el aire y es muy utilizado como detector en purificadores de aire. Internamente contiene un diodo LED y un fototransistor colocados de tal forma que son capaces de detectar las reflexiones de luz sobre las partículas de polvo.

Es un sensor que consume muy poco (entorno a 20mA) y puede ser alimentado hasta 7V DC. La salida es un voltaje lineal y proporcional a la cantidad de polvo detectado. La sensibilidad del sensor es de 0.5V por cada 0.1mg por metro cúbico, el sensor tiene un conector JST de 6 pines.


